\newglossaryentry{reinforcement learning}
{
    name={Reinforcement Learning},
    description={A learning paradigm in which an agent interacts with an environment and learns a policy to maximize expected cumulative reward through trial-and-error}
}

\newglossaryentry{agent}
{
    name={Agent},
    description={An entity that receives observations from its environment and performs actions to achieve a goal, typically defined as maximizing the expected cumulative return}
}

\newglossaryentry{policy gradient}
{
    name={Policy Gradient},
    description={A reinforcement learning method that directly optimizes a parametrized policy $\pi_{\theta}(a|s)$ by performing gradient ascent on the expected cumulative return. The policy parameters are updated in the direction that increases the probability of rewarding actions \cite{sutton2018reinforcement}}
}

\newglossaryentry{general-sum game}
{
    name={General-Sum Game},
    description={A type of strategic interaction in game theory where the total payoff to all players can very. Unlike in zero-sum-games one player's gain does not necessarily equal another's loss. The games allow for cooperative, competitive, win-win, or lose-lose outcomes}
}

\newglossaryentry{reward function}
{
    name={Reward Function},
    description={
    A function that specifies the expected immediate reward received after 
    taking action $A_t$ in state $S_t$. It is commonly denoted as 
    $r(s,a)$ or as the expected value of $R_{t+1}$ conditioned on 
    $(S_t, A_t)$ \cite{sutton2018reinforcement}
    }
}

\newglossaryentry{exploration-exploitation}
{
    name={Exploration vs. Exploitation},
    description={
    The fundamental trade-off in reinforcement learning between selecting 
    actions that are known to yield high rewards (exploitation) and selecting 
    actions to gather new information about the environment (exploration). 
    An effective learning strategy balances both in order to maximize 
    long-term return \cite{sutton2018reinforcement}
    }
}

\newglossaryentry{environment}
{
    name={Environment},
    description={
    The external system with which the agent interacts. 
    Given an action $A_t$, the environment transitions to a new state and 
    produces a reward according to the transition dynamics 
    $p(s', r \mid s,a)$ \cite{sutton2018reinforcement}
    }
}

\newglossaryentry{state}
{
    name={State},
    description={
    A representation of the environment at a given time step, denoted as $S_t$. 
    The state contains all information required to describe the decision-making 
    situation at time $t$. After the agent selects an action $A_t$, 
    the environment transitions to a new state $S_{t+1}$ according to 
    the transition dynamics of the MDP \cite{sutton2018reinforcement}
    }
}

\newglossaryentry{action}
{
    name={Action},
    description={An element $a_{t}\in\mathcal{A}$ selected by the agent at time step $t$. In this work, an action consists of a composite decision including collaboration choice, project selection and effort allocation}
}

\newglossaryentry{reward}
{
    name={Reward},
    description={
    A feedback signal provided by the environment after the agent 
    takes an action $A_t$ in state $S_t$. The reward, denoted as $R_{t+1}$, 
    indicates the immediate desirability of the resulting transition and 
    can be positive or negative \cite{sutton2018reinforcement}
    }
}

\newglossaryentry{policy}
{
    name={Policy},
    description={
    A mapping from states to actions that defines the behavior of an agent. 
    A policy can be deterministic or stochastic and is commonly denoted as 
    $\pi(a \mid s)$. In reinforcement learning, the objective is to improve 
    the policy in order to maximize the expected return. 
    Policies may be represented explicitly (e.g. as lookup tables) or 
    parameterized, for example by a neural network \cite{sutton2018reinforcement}
    }
}

\newglossaryentry{return}
{
    name={Return},
    description={
    The cumulative reward that an agent aims to maximize. 
    In the episodic case, the simplest form is the sum of future rewards:
    \[
    G_t \doteq R_{t+1} + R_{t+2} + \dots + R_T,
    \]
    where $T$ denotes the final time step \cite{sutton2018reinforcement}
    }
}

\newglossaryentry{markov decision process}
{
    name={Markov Decision Process},
    description={
    A formal problem for sequential decision-making under uncertainty. 
    An MDP is defined by a tuple $(\mathcal{S}, \mathcal{A}, P, R, \gamma)$, 
    where $\mathcal{S}$ is the state space, $\mathcal{A}$ the action space, $P(s', r \mid s,a)$ the transition dynamics and $\gamma$ the discount factor for future rewards. Formally:
    $$P(s', r|s,a)\doteq Pr\{S_{t}=s', R_t=r|S_{t-1}=s, A_{t-1}=a\}$$
    The Markov property states that the next state and reward depend only on the current state and action, not on the full history \cite{sutton2018reinforcement} 
    }
}

\newglossaryentry{discount_factor}
{
    name={Discount Factor},
    description={
    A scalar $\gamma \in [0,1]$ that determines how strongly future rewards are weighted relative to immediate rewards. 
    In the discounted return,
    \[
    G_t \doteq \sum_{k=0}^{\infty} \gamma^k R_{t+k+1},
    \]
    smaller values of $\gamma$ prioritize short-term rewards, while values close to 1 emphasize long-term returns \cite{sutton2018reinforcement}
    }
}

\newglossaryentry{on-off-policy learning}
{
    name={On-Policy / Off-Policy Learning},
    description={
    A distinction in reinforcement learning based on how data is collected. 
    In on-policy learning, the agent improves the same policy that is used 
    to generate behavior. In off-policy learning, the agent learns about a 
    target policy while following a different behavior policy 
    \cite{sutton2018reinforcement}
    }
}

\newglossaryentry{value function}
{
    name={Value Function},
    description={
    A function that estimates the expected return of a state or a state--action pair 
    under a given policy $\pi$. The state-value function is denoted as 
    $v_\pi(s)$ and represents the expected return when starting from state $s$ 
    and following policy $\pi$. The action-value function is denoted as 
    $q_\pi(s,a)$ and represents the expected return when taking action $a$ 
    in state $s$ and following policy $\pi$ 
    \cite{sutton2018reinforcement}
    }
}